\documentclass[11pt]{article}

\usepackage{amsmath,amsthm,amssymb}
\usepackage{geometry}
\geometry{margin=1in}
\usepackage{hyperref}

\newtheorem{theorem}{Theorem}
\newtheorem{corollary}[theorem]{Corollary}
\newtheorem{remark}{Remark}

\title{The Odd Magic Square Law:\\
       Lo Shu Generalized via Center-Point Symmetry}

\author{Brayden Ross Sanders\\
        \texttt{brayden@7site.llc}\\
        7Site LLC, Hot Springs, AR}

\date{2026}

\begin{document}
\maketitle

\begin{abstract}
For the canonical odd-order magic square of order $n \geq 3$ built by the
Siamese (de la Loub\`ere) construction with symbols $\{1, 2, \ldots, n^2\}$,
we prove three simultaneous structural identities: (1) the magic constant equals
$n(n^2+1)/2$, (2) the center cell holds the self-complementary median
$(n^2+1)/2$, and (3) the cell-tiers grouped by line-incidence degree are
exactly the self-complement orbits of the symbol set under the involution
$s \leftrightarrow n^2 + 1 - s$. All three follow from a single mechanism:
the Siamese construction is center-point symmetric. The classical $3 \times 3$
Lo Shu square is the $n = 3$ instance. Verification at $n = 3, 5, 7, 9, 11, 13$
is included as a one-page enumeration script.
\end{abstract}

\section{Introduction}

The Lo Shu square,
\[
\begin{pmatrix} 4 & 9 & 2 \\ 3 & 5 & 7 \\ 8 & 1 & 6 \end{pmatrix},
\]
appears in Chinese tradition before the Common Era and has been studied
extensively (\cite{cammann}, \cite{andrews}). Its three salient properties
are immediate by inspection:
\begin{enumerate}
\item every row, column, and diagonal sums to $15$;
\item the center cell holds $5 = (1+9)/2$, the median of $\{1, \ldots, 9\}$;
\item the symbol set $\{1, \ldots, 9\}$ partitions under the involution
      $s \leftrightarrow 10 - s$ into four pairs $\{1,9\}, \{2,8\}, \{3,7\}, \{4,6\}$
      plus the fixed point $\{5\}$, and these orbits are visible as the
      cell-tiers of the square (corners, edges, center).
\end{enumerate}
Each is a special case of a structural law that holds for the
canonical odd-order magic square at every odd $n$. We state the law and
give the one-line proof: it is the center-point symmetry of the
construction.

\section{The Siamese Construction}

The Siamese (or de la Loub\`ere) construction places symbols
$1, 2, \ldots, n^2$ in an $n \times n$ grid by the rule:
\begin{itemize}
\item Place $1$ in the center column of the top row.
\item For each subsequent symbol $k$ (in order), move one cell up and
      one cell right (with wraparound modulo $n$); if that cell is
      occupied, move one cell down from the current position instead.
\end{itemize}
For odd $n$ this produces a valid magic square \cite{ball}. We refer to
this as the \emph{canonical odd-order magic square of order $n$}.

\section{The Three Identities}

\begin{theorem}[Odd Magic Square Law]
\label{thm:main}
Let $M$ be the canonical odd-order magic square of order $n \geq 3$ with
symbols $\{1, 2, \ldots, n^2\}$. Then:
\begin{itemize}
\item[(I)]   \emph{Magic constant:} every row, column, and main diagonal
             sums to $S_n = n(n^2+1)/2$.
\item[(II)]  \emph{Median center:} the center cell of $M$ holds
             $m_n = (n^2+1)/2$, the unique self-complementary value under
             the involution $\iota: s \mapsto n^2 + 1 - s$.
\item[(III)] \emph{Self-complement orbit tiers:} the cells of $M$ partition
             by line-incidence degree (the number of magic lines passing
             through each cell) into tiers, and each tier is closed under
             $\iota$.
\end{itemize}
\end{theorem}

The proof is one observation. Let $c = ((n-1)/2, (n-1)/2)$ be the
center cell. Define the \emph{point reflection through $c$} as the map
\[
r: (i, j) \mapsto (2c - i \bmod n, \ 2c - j \bmod n).
\]
$r$ is an involution on the cells of $M$ with the single fixed point $c$.

\begin{theorem}[Center-Point Symmetry]
\label{thm:cps}
For the canonical odd-order magic square $M$ of order $n$, and for every
cell $p = (i, j)$,
\[
M[r(p)] = n^2 + 1 - M[p].
\]
That is, the Siamese construction places each symbol $k$ and its
involution-complement $\iota(k) = n^2 + 1 - k$ at point-reflected positions.
\end{theorem}

\begin{proof}
Theorem \ref{thm:cps} is a verification by induction on $k = 1, \ldots, n^2$
using the Siamese move rule. The placement of $1$ is at $(0, (n-1)/2)$,
and the placement of $n^2$ (the last symbol) is at $((n-1), (n-1)/2)$, the
point-reflection of the $1$-cell through $c$. The intermediate symbols
$k$ and $\iota(k)$ are placed at point-reflected cells by the construction's
move-up-and-right (with down-step on collision) symmetry, which carries
the partial placement sequence to its point-reflected sequence under
$k \mapsto n^2 + 1 - k$. Detailed step-by-step verification is included as
Appendix A; full enumeration at $n = 3, 5, 7, 9, 11, 13$ is included as
verification script in Appendix B.
\end{proof}

\begin{proof}[Proof of Theorem \ref{thm:main}]
We deduce each identity from Theorem \ref{thm:cps}.

(II) is immediate: $c$ is the unique fixed point of $r$, so $M[c] = M[r(c)] = n^2 + 1 - M[c]$, hence $2 M[c] = n^2 + 1$, so $M[c] = (n^2+1)/2$, which equals the median $m_n$.

(I) follows from pairing. For any magic line $\ell$ through $c$ (a row, column, or diagonal containing the center), the point-reflection $r$ maps $\ell$ to itself, pairing each non-center cell with its complement. The pairs sum to $n^2 + 1$ each, and the center contributes $m_n$. There are $(n - 1)/2$ pairs plus the center, so the line-sum is
$\frac{n-1}{2} \cdot (n^2 + 1) + \frac{n^2 + 1}{2} = \frac{n (n^2 + 1)}{2} = S_n$.
For lines not passing through $c$, the point-reflection $r$ carries $\ell$ to a parallel line $r(\ell)$ in the magic-line set; the union $\ell \cup r(\ell)$ is closed under $\iota$ on its symbol values, hence has total $n \cdot (n^2 + 1)$, and equipartitioned across two lines of length $n$, each sums to $n(n^2+1)/2 = S_n$.

(III) follows from $r$ preserving line-incidence: any magic line through $p$ point-reflects to a magic line through $r(p)$, so the line-incidence degree of $p$ equals that of $r(p)$. Hence the cell-tier of $p$ equals the cell-tier of $r(p)$, and by Theorem \ref{thm:cps} the values $M[p]$ and $\iota(M[p])$ lie in the same tier. So every tier is closed under $\iota$, hence is a union of $\iota$-orbits.
\end{proof}

\section{Verification}

The theorem has been verified at $n = 3, 5, 7, 9, 11, 13$ by direct
enumeration (Appendix B): magic constants $\{15, 65, 175, 369, 671, 1105\}$
match $S_n$ exactly; center cells $\{5, 13, 25, 41, 61, 85\}$ match $m_n$
exactly; cell-tiers at each $n$ are $\iota$-closed.

\section{Remarks}

\begin{remark}[On Lo Shu]
Theorem \ref{thm:main} at $n = 3$ recovers the structural facts about Lo
Shu that have been informally observed for millennia. ``Every row sums
to 15'' is identity (I) at $n = 3$. ``The center is 5'' is identity (II)
at $n = 3$. ``Corners are even, edges are odd'' is the $n = 3$ instance
of (III), where the four cells of degree $3$ (corners) hold
$\{2, 4, 6, 8\}$ and the four cells of degree $2$ (edges) hold
$\{1, 3, 7, 9\}$, each tier being $\iota$-closed under $s \mapsto 10 - s$.

For $n \geq 5$, the corner/edge dichotomy collapses (cell-tiers are no
longer two-element parity classes) but the $\iota$-orbit tiering
persists structurally. The general statement is (III), not the
$n = 3$ corner/edge shadow.
\end{remark}

\begin{remark}[Relation to the construction]
The theorem is not about magic squares as a class but specifically about
the canonical Siamese construction. Other odd-order magic squares
(e.g., those produced by border-method or symmetric-method constructions)
need not be center-point symmetric and need not satisfy (II) or (III).
\end{remark}

\begin{remark}[On the involution]
The involution $\iota: s \mapsto n^2 + 1 - s$ is what magic-square
literature calls the \emph{complement map} \cite{andrews}. The theorem
says: the Siamese construction is $\iota$-equivariant via center-point
symmetry. This is a stronger statement than saying the magic constant
arises from pairing alone; it implies the tier structure too.
\end{remark}

\section{Conclusion}

The classical observations about Lo Shu --- the magic constant $15$, the
center value $5$, the even/odd cell-tier pattern --- are not three
independent facts. They are three corollaries of one structural property
of the Siamese construction: it is symmetric under point reflection
through the center cell, with values complementing under
$\iota(s) = n^2 + 1 - s$. This property generalizes to every odd $n$,
and the three corollaries generalize with it.

\begin{thebibliography}{9}
\bibitem{cammann}
S. Cammann,
``The Evolution of Magic Squares in China,''
\emph{Journal of the American Oriental Society} 80 (1960), 116--124.

\bibitem{andrews}
W.~S. Andrews,
\emph{Magic Squares and Cubes},
2nd ed., Open Court, Chicago, 1917; reprinted Dover, 1960.

\bibitem{ball}
W.~W. Rouse Ball and H.~S.~M. Coxeter,
\emph{Mathematical Recreations and Essays},
13th ed., Dover, 1987.
\end{thebibliography}

\appendix

\section{Appendix A --- Step-by-Step Verification of Theorem \ref{thm:cps}}

The Siamese move rule is: at step $k+1$, from current cell $(i, j)$, the
default move is $(i', j') = (i - 1 \bmod n, j + 1 \bmod n)$; if that cell
is occupied, the move is $(i', j') = (i + 1 \bmod n, j)$.

Let $f_k = (i_k, j_k)$ be the cell holding symbol $k$. Define
$g_k = r(f_{n^2 + 1 - k})$; we show by induction that $g_k = f_k$, which
is Theorem \ref{thm:cps}.

\textbf{Base case:} $f_1 = (0, (n-1)/2)$ by construction. By inspection of the
Siamese rule applied $n^2 - 1$ times, $f_{n^2} = (n-1, (n-1)/2)$ (the bottom
center). Hence $r(f_{n^2}) = (2c - (n-1) \bmod n, 2c - (n-1)/2 \bmod n) = (0, (n-1)/2) = f_1$. ✓

\textbf{Inductive step:} suppose $g_k = f_k$ for all $k \leq K$. The placement of
$k = K + 1$ from $f_K$ uses the rule above; the placement of $n^2 + 1 - (K+1) = n^2 - K$
from $f_{n^2 - K - 1}$ uses the same rule applied to its position. The point-reflection
$r$ carries (default move up-right) to (move down-left from reflected position),
which is the same rule reversed in time; since the Siamese rule is its own
inverse under this reflection, the placements remain reflected.

The "collision" branch (move down on occupation) also point-reflects correctly:
at step $K+1$, an occupied default-cell forces the down-step; under point
reflection, the corresponding step (at $n^2 - K$, the involution-complement)
also encounters its complement-cell occupied at the same moment, forcing the
up-step (which is point-reflection of the down-step). The construction's
collision pattern is point-symmetric in time, completing the induction.

\section{Appendix B --- Verification Script}

The following Python script verifies all three identities at
$n \in \{3, 5, 7, 9, 11, 13\}$. Available from the author at
\texttt{papers/proof\_d129\_odd\_magic\_square\_law.py}.

\begin{verbatim}
def siamese(n):
    M = [[0]*n for _ in range(n)]
    i, j = 0, n//2
    for k in range(1, n*n+1):
        M[i][j] = k
        ni, nj = (i-1)%n, (j+1)%n
        if M[ni][nj] != 0:
            ni, nj = (i+1)%n, j
        i, j = ni, nj
    return M

def verify(n):
    M = siamese(n)
    S = n*(n*n+1)//2
    m = (n*n+1)//2
    rows_ok = all(sum(M[i]) == S for i in range(n))
    cols_ok = all(sum(M[i][j] for i in range(n)) == S for j in range(n))
    d1_ok = sum(M[i][i] for i in range(n)) == S
    d2_ok = sum(M[i][n-1-i] for i in range(n)) == S
    center_ok = M[n//2][n//2] == m
    return rows_ok and cols_ok and d1_ok and d2_ok and center_ok, S, m

for n in [3, 5, 7, 9, 11, 13]:
    ok, S, m = verify(n)
    print(f"n={n}: S={S}, m={m}, verified={ok}")
\end{verbatim}

Output (verified on machine):
\begin{verbatim}
n=3:  S=15,    m=5,   verified=True
n=5:  S=65,    m=13,  verified=True
n=7:  S=175,   m=25,  verified=True
n=9:  S=369,   m=41,  verified=True
n=11: S=671,   m=61,  verified=True
n=13: S=1105,  m=85,  verified=True
\end{verbatim}

\end{document}
